% Tapis de Quillen, with Quillen's offprint and preprint (cote 162-5, typed part): archive cote 162-5, site page 32 (PDF page 33).
% typescript. Draft compared: none.
% AI-transcribed by Claude Opus 5.5 (claude-opus-5-5, adaptive thinking, effort high), Message Batches API, batch t162_full (msgbatch_01UAh3zyWPd454AeDkhnmQ33), 2026-11-13; not yet checked by eye.
\scanpage{162-5/32}
\margin{notes du 10.9.68}

Tapis de Quillen.

I \quad \uline{Relations entre catégories et ensembles \unsure{semi-simplicaux}.}

A toute catégorie C,  on associe un ensemble semi-simplicial S(C),
trouvant ainsi un foncteur pleinement fidèle \xout{dexCat}
\[ S; \ (\mathrm{Cat}) \longrightarrow (\mathrm{Ssimpl}) \quad . \]
\xout{Si C est associé à} Les systêmes locaux d'ensemble sur SC correspondent
aux \unsure{fancteurs|foncteurs} sur C qui transforment toute flêche en isomorphisme (i.e.
qui se factorisent par le groupoïde associé à  C). Les $H^i$ sur SC d'un tel
système local ($H^0$ pour ens, $H^1$ pour groupes, $H^i$ quelconques pour groupes
abéliens) s'interprètent en termes des foncteurs $\underleftarrow{\mathrm{lim}}^{(i)}$ \unsure{dérivés|dérigés} de $\underleftarrow{\mathrm{lim}}$ ,
ou si on préfère, des $H^i($ du \uline{topos} C . On voit ainsi à quelle condition
un foncteur  $C \to C'$ induit un\xout{\unsure{e}}  homotopisme  $SC \to SC'$ :  \xout{\unsure{e}}en vertu du
critère cohomologique de Artin-Mazur, il f et s que  pour tout système
de coéfficients  F'sur C', l'homomorphisme naturel $\underset{C'}{\underleftarrow{\mathrm{lim}}}^{(i)}F' \longrightarrow \underset{C}{\underleftarrow{\mathrm{lim}}}^{(i\ )}F$
soit un isomorphisme (pour les i pour lesquels cela a un sens). \xout{\unsure{En}\illegible{1}}

A C on peut associer le topos $\overset{\hand{\sim}}{C}$ , qui varie de façon \uline{covariante} avec C.
(NB le foncteur $C \mapsto \overset{\hand{\sim}}{C}$  n'a plus rien de pleinement fidèle, semble-til ??)
Les systêmes de coéfficients \inserted[typed]{ensemblistes} sur C ( = \inserted[typed]{\unsure{déf}} les foncteurs \inserted[typed]{$C^o \longrightarrow \mathrm{Ens}$} transformant isomor-%
phismes en isomorphismes) correspondent aux faisceaux localement constants,
i.e. les objets localement constants de $\overset{\hand{\sim}}{C}$, définis intrinsèquement en termes
de $\overset{\hand{\sim}}{C}$ .  Ainsi, le fait pour un foncteur $F:C \to C'$ d'induire un\xout{e équivalence}
homotopisme $S(C) \to S(C')$  ne  dépend que du morphisme de topos $\overset{\hand{\sim}}{F}: \overset{\hand{\sim}}{C} \to \overset{\hand{\sim}}{C'}$
induit, et signifie que pour tout faisceau localement constant F' sur C' i.
e. sur $\overset{\hand{\sim}}{C'}$ ,  les applications induites  $H^i(\overset{\hand{\sim}}{C'},F') \longrightarrow H^i(\overset{\hand{\sim}}{C},\overset{\hand{\sim}}{F}{}^{\hand{*}}(F'))$
sont des isomorphismes (pour les i pour lesquels cela a un\xout{xx} sens).

\mspage{32}
